\documentstyle[%


righttag%


,11pt%


,amssymb%


,russian%


,izvru%


]{amsart}





\newcommand{\tl}[3]{\medskip\noindent{\bf#1}\ #2 \dotfill #3 \par}


\pagestyle{empty}


\begin{document}


%\tableofcontents


\par


\vskip 2cm


\par


\bigskip


\bigskip





\par


\centerline{\Large СОДЕРЖАНИЕ}


\bigskip








\tl{Бекбаев У.Д.} 


{Об эквивалентности и инвариантах линейных  


дифференциальных\newline уравнений в частных производных  


второго порядка от двух переменных относительно\newline


замены переменных}{3}


\tl{Бильченко Н.Г., Гараев К.Г., Овчинников В.А.}


{Группа симметрии и законы\newline сохранения в задаче оптимального 


управления ламинарным пограничным слоем\newline неравновесно 


диссоциирующего газа}{11}


\tl{Гаврилюк М.Н.} 


{О покрытии кругов в классе Бибербаха--Эйленберга}{20}


\tl{Корнеев В.Г., Енсен С.}


{Эффективное предобуславливание методом декомпозиции \newline


области для $p$-версии с иерархическим базисом. II}{24}


\tl{Крепкогорский В.Л.}


{О многомерных методах интерполяции}{41}


\tl{Мокейчев В.С., \fbox{Мокейчев А.В.}}   


{Новый подход к теории линейных задач


для систем\newline дифференциальных уравнений


в частных производных. III}{50}


\tl{Окрут С.И.}


{Голоморфные конформные субмерсии}{60}


\tl{Сабитов К.Б., Карамова А.А., Шарафутдинова Г.Г.}


{К теории уравнений смешанного\newline типа с двумя линиями вырождения}{70}








\bigskip


{\centerline{КРАТКИЕ СООБЩЕНИЯ}}


\bigskip





\tl{Алеев Р.Ж., Панина Г.А.} 


{Единицы циклических групп порядков 7 и 9}{81}


\bigskip





\noindent Математическая жизнь в России и зарубежом.








\vfill


\noindent\raisebox{2pt}{\copyright} Казанский госудаpственный унивеpситет


\end{document}


